\begin{theorem}[Closed norm-trace correction formula]\label{mod:cases-wildodd-correctionformula}
Let $d$ be wild odd correction data, put
$p=\operatorname{char}k_F$, $n=N_{K/F}(u)$, and write
$s_r=s_r(u)$ for the elementary symmetric functions of the conjugates of
$u$, so that $s_1=\operatorname{Tr}_{K/F}(u)$ and $s_p=n$.  Lean defines the
manuscript numerator by the literal finite sum
\[
 B=s_1+\sum_{q\in[1,p-1)}(-n)^q s_{p-q}.
\]
It also proves an explicit finite-sum reindexing from the shifted range
$q=0,\ldots,p-3$ produced by the prime-field calculation to this interval:
\[
 s_1+\sum_{q=0}^{p-3}(-n)^{q+1}s_{p-(q+1)}=B.
\]

The exceptional index is kept separate from every nonzero prime-field
index.  For $j\in\mathbf F_p^\times$, with $[j]$ the canonical
Teichmüller lift, Lean first derives from
$N(u+[j])=(n+[j])z_j$ the exact pointwise identity
\[
 x_j=
 \frac{\sum_{r=1}^{p-1}[j]^{p-r}s_r}{n+[j]}.
\]
Every denominator $n+[j]$ is proved nonzero before this division.  The
finite geometric identity
\[
 \frac1{n+[j]}=
 \frac{\sum_{q=0}^{p-2}(-n)^q[j]^{p-2-q}}
      {1-n^{p-1}}
\]
is likewise used only after proving both $n+[j]\ne0$ and
$1-n^{p-1}\ne0$.  Applying the prime-field power-sum theorem to the
Teichmüller unit homomorphism leaves exactly the cases $r+q=1$ and
$r+q=p$, and gives the nonzero-index component
\[
 \sum_{j\in\mathbf F_p^\times}[j]x_j
   =\frac{(p-1)B}{1-n^{p-1}}.
\]

For the exceptional index, Lean uses the ratio in its proved orientation
\[
 N(n-u)=(n^p-n)z_0
\]
and the exact symmetric-polynomial expansion of $N(n-u)$.  After a
separate reflected-interval reindexing, it obtains
\[
 nx_0=\frac{B}{1-n^{p-1}}-s_1.
\]
Thus the correction quantity is defined exactly as
\[
 X=s_1+nx_0+
   \sum_{j\in\mathbf F_p^\times}[j]x_j,
\]
and Lean proves the field equality
\[
 \boxed{X=\frac{pB}{1-n^{p-1}}}.
\]
This is an equality in $F$, not a congruence.  No valuation or depth
estimate, and no result from CorrectionValuation or any later wild-odd
node, is used.
\LeanFile{LanglandsFirstMainLemma/Cases/WildOdd/CorrectionFormula.lean}
\lean{LanglandsFirstMainLemma.correctionFormulaB}
\lean{LanglandsFirstMainLemma.correctionFormula_B_reindex}
\lean{LanglandsFirstMainLemma.correctionFormula_n_mul_xZero_eq_B_sub_trace}
\lean{LanglandsFirstMainLemma.correctionFormula_nonzero_sum}
\lean{LanglandsFirstMainLemma.wildOddCorrectionX}
\lean{LanglandsFirstMainLemma.wildOdd_correctionFormula}
\leanok
\SourceText{Propositions prop:exact-odd-assembly and prop:X-closed, including equations eq:X-def, eq:B-def, eq:inverse-nj, eq:sum-jxj, and eq:X-closed.}
\uses{mod:cases-wildodd-correctionunits,mod:basic-finiteproducts}
\FormalizationContract The Lean definition of $X$ preserves the exceptional
index $0$ and the nonzero Teichmüller indices as distinct summands.  The
proof expands the already-defined correction units through the exact norm
identities, establishes every displayed denominator as nonzero before
division or cancellation, performs the prime-field reindexing explicitly,
and concludes with the exact rational identity required by the following
CorrectionValuation node.  It does not import or assume that node.
\end{theorem}
